\section{Multiple pre-Lie algebras}

We here fix a nonempty set $\T$ of edges types.

\subsection{Definition}

\begin{Definition} \label{defi3}
A $\T$-multiple pre-Lie algebra is a family $(V,(\bullet_t)_{t\in \T})$,
where $V$ is a vector space and for all $t\in \T$,
$\bullet_t$ is a bilinear product on $V$ such that
\begin{gather*}
\forall t,t'\in \T,\quad\forall x,y,z\in V,\qquad
x\bullet_{t'}(y \bullet_t z)-(x\bullet_{t'} y)\bullet_t z=x\bullet_t(z \bullet_{t'} y)-(x\bullet_t z)\bullet_{t'} y.
\end{gather*}
\end{Definition}

\begin{Remark}
For any $t\in \T$, $(V,\bullet_t)$ is a pre-Lie algebra.
More generally, for any family $\lambda=(\lambda_t)_{t\in \T}\in \K^{(\T)}$,
putting $\bullet_\lambda=\sum \lambda_t \bullet_t$, $(V,\bullet_\lambda)$ is a pre-Lie algebra.
\end{Remark}

\begin{Proposition}
Let $\D$ be any set; we denote by $\g_{\D,\T}$ the vector space generated by $\bfT_{\D,\T}$.
For any $T,T'\in \bfT_{\D,\T}$,
$v\in V(T)$ and $t\in \T$, we denote by $T\bullet_t^{(v)} T'$ the $\D$-decorated $\T$-typed tree
obtained by grafting $T'$ on $v$ $($that is to say adding an edge between $v$ and the root of $T')$, the created edge being of type $t$.
We then define a product $\bullet_t$ on $\g_{\D,\T}$ by
\begin{gather*}
\forall T,T'\in \bfT_{\D,\T},\qquad
T\bullet_t T'=\sum_{v\in V(T)} T\bullet_t^{(v)} T'.
\end{gather*}
Then $(\g_{\D,\T},(\bullet_t)_{t\in \T})$ is a $\T$-multiple pre-Lie algebra.
\end{Proposition}

\begin{proof} Let $T$, $T'$, $T''$ be elements of $\bfT_{\D,\T}$ and $t',t''\in \T$.
\begin{gather*}
(T\bullet_{t'} T')\bullet_{t''} T''-T\bullet_{t'}(T'\bullet_{t''}T'')
\\ \qquad
{}=\sum_{v\in V(T), v'\in V(T)\sqcup V(T')} \big(T\bullet_{t'}^{(v)} T'\big)\bullet_{t''}^{(v')} T''
-\sum_{v\in V(T),v'\in V(T')}T\bullet_{t'}^{(v)}\big(T'\bullet_{t''}^{(v')} T''\big)
\\ \qquad
{}=\sum_{v\in V(T), v'\in V(T)}\!\!\!\!\!\! \big(T\bullet_{t'}^{(v)} T'\big)\bullet_{t''}^{(v')} T''
\!+\!\!\!\sum_{v\in V(T),v'\in V(T')}\!\!\!\!\!\!\big(T\bullet_{t'}^{(v)} T'\big)\bullet_{t''}^{(v')} T''
\!-T\bullet_{t'}^{(v)}\!\big(T'\bullet_{t''}^{(v')} T''\big)
\\ \qquad
{}=\sum_{v\in V(T), v'\in V(T)} \big(T\bullet_{t'}^{(v)} T'\big)\bullet_{t''}^{(v')} T''
\\ \qquad
{}=\sum_{v\in V(T), v'\in V(T)} \big(T\bullet_{t''}^{(v')} T''\big)\bullet_{t'}^{(v)} T'
\\ \qquad
{}=(T\bullet_{t''} T'')\bullet_{t'} T'-T\bullet_{t''}(T''\bullet_{t'}T').
\end{gather*}
So $\g_{\D,\T}$ is indeed a $\T$-multiple pre-Lie algebra.
\end{proof}

\begin{Example} If $a,b,c\in \D$ and $\typ{10}$\!, $\typ{2} \in \T$,\vspace{-1ex}
\begin{gather*}
\tddeux{$a$}{$b$}{10}\bullet_{\typ{10}} \tdun{$c$}=
\tdtroisun{$a$}{$c$}{$b$}{10}{10}+\tdtroisdeux{$a$}{$b$}{$c$}{10}{10},\qquad\qquad
\tddeux{$a$}{$b$}{10}\bullet_{\typ{2}} \tdun{$c$}=
\tdtroisun{$a$}{$c$}{$b$}{10}{2}+\tdtroisdeux{$a$}{$b$}{$c$}{10}{2}.
\end{gather*}\end{Example}


\subsection{Guin--Oudom extension of the pre-Lie products}

{\samepage\begin{Notation}
Let $(\partial_t)_{t\in \T}$ be a family of formal symbols indexed by $\T$. For any vector space~$V$, we denote
\begin{gather*}
V^{\oplus \T}=V\otimes \mathrm{Vect}(\partial_t,\:t\in \T).
\end{gather*}
Then
\begin{gather*}
V^{\oplus \T}=\bigoplus_{t\in \T} V\otimes \partial_t.
\end{gather*}
\end{Notation}
In order to enlighten the notations, we write $v\partial_t$ instead of $v\otimes \partial_t$
for any $v\in V$ and for any $t\in \T$.

}

\begin{Lemma}
If for any $t\in \T$, $\bullet_t$ is a bilinear product on a vector space $V$,
we define $\bullet\,{:}$ $\big(V^{\oplus \T}\big)^{\otimes 2}\longrightarrow V^{\oplus \T}$ by
\begin{gather*}
x\partial_t \bullet x'\partial_{t'}=(x\bullet_{t'} y)\partial_t.
\end{gather*}
Then $(V,(\bullet_t)_{t\in \T})$ is a $\T$-multiple pre-Lie algebra if,
and only if, $\big(V^{\oplus \T},\bullet\big)$ is a pre-Lie algebra.
\end{Lemma}

\begin{proof}
Let $x,x',x''\in V$, $t,t',t''\in \T$. Then
\begin{gather*}
x\partial_t\bullet (x'\partial_{t'}\bullet x''\partial_{t''})-(x\partial_t\bullet x'\partial_{t'})\bullet x''\partial_{t''}
=\left((x\bullet_{t'} x')\bullet_{t''} x''-x\bullet_{t'}(x'\bullet_{t''}x'')\right)\partial_t,
\end{gather*}
which implies the result.
\end{proof}

\begin{Notation}
The symmetric algebra $S(V)$ is given with its usual coproduct $\Delta$, making it a~bialgebra:
\begin{gather*}
\forall x\in V,\qquad \Delta(x)=x\otimes 1+1\otimes x.
\end{gather*}
We shall use Sweedler's notation: for any $w\in S(V)$, $\Delta(w)=\sum w^{(1)}\otimes w^{(2)}$.
The counit of this coproduct is denoted by $\varepsilon$: this is the unique algebra morphism from $S(V)$ to $\K$ sending any $v\in V$ to $0$.
\end{Notation}

\begin{Theorem}
Let $V$ be a $\T$-multiple pre-Lie algebra. One can define a product
\begin{gather*}
\bullet\colon\ S(V)\otimes S\big(V^{\oplus\T}\big)\longrightarrow S(V)
\end{gather*}
in the following way: for any $u,v \in S(V)$, $w\in S\big(V^{\oplus\T}\big)$, $x\in V$, $t\in \T$,
\begin{gather*}
1\bullet w=\varepsilon(w),
\\
u\bullet 1=u,
\\
uv\bullet w=\sum\big(u\bullet w^{(1)}\big)\big(v\bullet w^{(2)}\big),
\\
u\bullet w (x\partial_t)=(u\bullet w)\bullet_t x-x\bullet (w\bullet_t x),
\end{gather*}
where $\bullet_t$ is extended to $S(V)\otimes V$ and $S\big(V^{\oplus\T}\big)\otimes V$
by the following: for any $x_1,\dots,x_k,x\in V$, for any $t_1,\dots,t_k \in \T$,
\begin{gather*}
x_1\cdots x_k \bullet_t x=\sum_{i=1}^k x_1\cdots (x_i\bullet_t x)\cdots x_k,
\\
(x_1\partial_{t_1})\cdots (x_k\partial_{t_k}) \bullet_t x
=\sum_{i=1}^k (x_1\partial_{t_1})\cdots ((x_i\bullet_t x)\partial_{t_i})\cdots (x_k\partial_{t_k}).
\end{gather*}
\end{Theorem}

\begin{proof} \textit{Uniqueness}. The last formula allows to compute $x\bullet w$ for any $x\in V$
and $w\in S\big(V^{\oplus\T}\big)$ by induction on the length of $w$;
the other ones allow to compute $u\bullet w$ for any $u\in S(V)$ by induction on the length on $u$.
So this product $\bullet$ is unique.

\textit{Existence}. Let us use the Guin--Oudom construction \cite{Guin1,Guin2} on the pre-Lie algebra $V^{\otimes \T}$.
We obtain a product $\bullet$
defined on $S\big(V^{\oplus\T}\big)$ such that for any $u,v,w\in S\big(V^{\oplus\T}\big)$, $x\in V^{\oplus \T}$:
\begin{gather*}
1\bullet w=\varepsilon(w),
\\
u\bullet 1=u,
\\
uv\bullet w=\sum \big(u\bullet w^{(1)}\big)\big(v\bullet w^{(2)}\big),
\\
u\bullet w x=(u\bullet w)\bullet x-x\bullet (w\bullet x).
\end{gather*}
Let $f\colon \T\longrightarrow \K$ be any nonzero map. We consider the surjective algebra morphism
$F\colon S\big(V^{\oplus\T}\big)\allowbreak\longrightarrow S(V)$,
sending $x\partial_t$ to $f(t)x$ for any $x\in V$, $t\in \T$. Its kernel is generated by the elements
$X_{t,t'}x=(f(t')\partial_t-f(t)\partial_{t'})x$,
where $x\in V$ and $t,t'\in \T$. We denote by $J$ the vector space generated by the elements $X_{t,t'}x$.
Let us prove that for any $w\in S\big(V^{\oplus\T}\big)$, $J\bullet w\subseteq J$ by induction on the length $n$ of $w$.
If $n=0$, we can assume that $w=1$ and this is obvious. If $n=1$, we can assume that $w=x'\partial_{t''}$. Then
\begin{gather*}
X_{t,t'}x\bullet w=(f(t')\partial_t-f(t)\partial_{t'}) x\bullet_{t''} x'
=X_{t,'t'}x\bullet_{t''} x' \in J.'
\end{gather*}
Let us assume the result at rank $n-1$. We can assume that $w=w' x'\partial_t$, the length of $w'$ being $n-1$.
For any $x\in J$,
\begin{gather*}
x\bullet w=(x\bullet w')\bullet x'-x\bullet(w'\bullet x').
\end{gather*}
The length of $w'$ and $w'\bullet x'$ is $n-1$, so $x\bullet w'$ and $x\bullet (w'\bullet x')$ belong to $J$.
From the case $n=1$,
$(x\bullet w')\bullet x'\in J$, so $x\bullet w\in J$.

For any $x\in J$, $u,v\in S\big(V^{\oplus\T}\big)$,
\begin{gather*}
xu\bullet v=\underbrace{\big(x\bullet v^{(1)}\big)}_{\in J} \big(u\bullet v^{(2)}\big) \in \ker(F).
\end{gather*}
This proves that $\ker(F)\bullet S\big(V^{\oplus\T}\big)\subseteq \ker(F)$. Hence, $\bullet$ induces a product also denoted by $\bullet$,
defined from $S(V)\otimes S\big(V^{\otimes \T}\big)$ to $S(V)$. It is not difficult to show that it does not depend on the choice of $f$
and satisfies the required properties.
\end{proof}

\begin{Definition}
Let $d\in \D$, $T_1,\dots,T_k\in \bfT_{\D,\T}$, $t_1,\dots,t_k\in \T$. We denote by
\begin{gather*}
B_d\bigg(\prod_{i\in [k]} T_i\partial_{t_i}\bigg)
\end{gather*}
the $\T$-typed $\D$-decorated tree obtained by grafting $T_1,\dots,T_k$ on a common root decorated by~$d$,
the edge between this root and the root of $T_i$ being of type $t_i$ for any~$i$.
This defines a map $B_d\colon S\left(\mathrm{Vect}(\bfT_{\D,\T})^{\oplus \T}\right)\longrightarrow S(\mathrm{Vect}(\bfT_{\D,\T}))$.
\end{Definition}

\begin{Lemma} \label{lemma8}
For any $d\in \D$, $T_1,\dots,T_k\in \bfT_{\D,\T}$, $t_1,\dots,t_k\in \T$,
\begin{gather*}
B_d\bigg(\prod_{i\in [k]} T_i\partial_{t_i}\bigg)=\tdun{$d$}\bullet\prod_{i\in [k]} T_i\partial_{t_i}.
\end{gather*}
\end{Lemma}

\begin{proof} We write $F=\prod_{i\in [k]} T_i\partial_{t_i}$: the integer $k$ is unique
and denoted by $\ell(F)$. We proceed by induction on $\ell(F)$.
If $\ell(F)=0$, then $F=1$ and $\tdun{$d$}\bullet 1=\tdun{$d$}=B_d(1)$.
let us assume the result for any forest $F'$ with $\ell(F')=\ell(F)-1$. Putting $k=\ell(F)$,
we can write $F=F'T\partial_t$, with $\ell(F')=\ell(F)-1$, $T=T_k$ and $t=t_k$. Then
\begin{align*}
\tdun{$d$}\bullet F&=(\tdun{$d$}\bullet F')\bullet T\partial_t-\tdun{$d$}\bullet (F'\bullet T\partial_t)\\
&=B_d(F')\bullet_t T-B_d(F'\bullet_t T)\\
&=B_d(F'T\partial_t)+B_d(F'\bullet_t T)-B_d(F'\bullet_t T)\\
&=B_d(F).
\end{align*}
So the result holds for all forests. \end{proof}

\begin{Corollary} \label{cor9}
Let $A$ be a $\T$-multiple pre-Lie algebra and, for any $d\in \D$, $a_d \in A$.
There exists a unique $\T$-multiple algebra morphism
$\phi\colon \g_{\D,\T}\longrightarrow A$, such that for any $d\in \D$, $\phi(\tdun{$d$})=a_d$.
In~other words, $\g_{\T,\D}$ is the free $\T$-multiple pre-Lie algebra generated by $\D$.
\end{Corollary}

\begin{proof} \textit{Uniqueness}. Using the Guin--Oudom product and lemma \ref{lemma8},
$\phi$ is the unique linear map inductively defined by
\begin{gather*}
\phi\bigg(B_d\bigg(\prod_{i\in[k]}T_i \partial_{t_i}\bigg)\bigg)
=a_d \bullet \prod_{i\in [k]} \phi(T_i) \partial_{t_i}.
\end{gather*}

\textit{Existence}. Let $T,T'\in \bfT_{\D,\T}$ and $t\in \T$. Let us prove that $\phi(T\bullet_t T')=\phi(T)\bullet_t \phi(T')$
by induction on $n=|T|$. If $n=1$, we assume that $T=\tdun{$d$}$. Then $T\bullet_t T'=B_d(T'\partial_t)$, so
\begin{gather*}
\phi(T\bullet_t T')=a_d\bullet (\phi(T')\partial_t=a_d \bullet_t \phi(T')=\phi(T)\bullet_t \phi(T').
\end{gather*}
Let us assume the result at all ranks $<|T|$. We put
\begin{gather*}
T=B_d\bigg(\prod_{i=1}^k T_i\partial_{t_i}\bigg).
\end{gather*}
By definition of the pre-Lie product of $\g_{\D,\T}$ in terms of grafting,
\begin{gather*}
T\bullet T'=B_d\bigg(\prod_{i=1}^k T_i \partial_{t_i} T'\partial_t\bigg)+\sum_{j=1}^k B_d\bigg(\prod_{i\neq j} T_i\partial_{t_i} (T_j\bullet_t T')\partial_{t_j}\bigg),
\\
\phi(T\bullet T')=a_d\bullet \prod_{i=1}^k \phi(T_i)\partial_{t_i}\phi(T')\partial_t
+\sum_{j=1}^k a_d\bullet \prod_{i\neq j} \phi(T_i)\partial_{t_i} (\phi(T_j\bullet_t T'))\partial_{t_j}
\\ \hphantom{\phi(T\bullet T')}
{}=a_d\bullet \prod_{i=1}^k \phi(T_i)\partial_{t_i}\phi(T')\partial_t
+\sum_{j=1}^k a_d\bullet \prod_{i\neq j} \phi(T_i)\partial_{t_i} (\phi(T_j)\bullet_t \phi(T'))\partial_{t_j}
\\ \hphantom{\phi(T\bullet T')}
{}=a_d\bullet \prod_{i=1}^k \phi(T_i)\partial_{t_i}\phi(T')\partial_t+a_d\bullet\bigg(\bigg(\prod_{i=1}^k
\phi(T_i)\partial_{t_i}\bigg)\bullet \phi(T')\partial_t\bigg)
\\ \hphantom{\phi(T\bullet T')}
{}=\bigg(a_d\bullet\prod_{i=1}^k \phi(T_i)\partial_{t_i} \bigg)\bullet \phi(T')\partial_t
\\ \hphantom{\phi(T\bullet T')}
{}=\phi(T)\bullet_t \phi(T').
\end{gather*}
So $\phi$ is a $\T$-multiple pre-Lie algebra morphism. \end{proof}

\subsection{Operad of typed trees}

We now describe an operad of typed trees, in the category of species.
We refer to \cite{Dotsenko,Vallette,Markl} for notations and definitions on operads.

\begin{Notation} Let $A$ be a finite set. If $T\in \bfT_\T(A)$ and $a\in T$:
\begin{enumerate}\itemsep=0pt
\item The subtrees formed by the connected components of the set of vertices, descendants of~$a$ ($a$ excluded)
are denoted by $T_1^{(a)},\dots,T_{n_a}^{(a)}$.
The type of the edge from $a$ to the root of~$T_i^{(a)}$ is denoted by $t_i$.
\item The tree formed by the vertices of $T$ which are not in $T_1^{(a)},\dots,T_{n_a}^{(a)}$,
at the exception of~$a$, is denoted by $T_0^{(a)}$.
\end{enumerate}
\end{Notation}

\begin{Example}
Let us consider the following tree:\vspace{-1ex}
\begin{gather*}T
=\begin{tikzpicture}[line cap=round,line join=round,>=triangle 45,x=0.7cm,y=0.7cm]
\clip(0.8,1.9) rectangle (4.,5.);
\draw [line width=0.5pt] (2.,2.)-- (1.7,2.5);
\draw [line width=0.5pt, dash pattern=on 2pt off 2pt] (2.,2.)-- (2.3,2.5);
\draw [line width=0.5pt] (1.7,2.5)-- (1.7,3.);
\draw [line width=0.5pt] (1.7,3.)-- (1.5,3.5);
\draw [line width=0.5pt, dash pattern=on 2pt off 2pt] (1.7,3.)-- (2.,3.5);
\draw [line width=0.5pt, dash pattern=on 2pt off 2pt] (1.7,3.)-- (2.5,3.5);
\draw [line width=0.5pt] (1.7,3.)-- (3.,3.5);
\draw [line width=0.5pt] (2.,3.5)-- (2.,4.);
\draw [line width=0.5pt, dash pattern=on 2pt off 2pt] (2.5,3.5)-- (2.5,4.);
\draw [line width=0.5pt] (3.,3.5)-- (3.,4.);
\draw [line width=0.5pt, dash pattern=on 2pt off 2pt] (3.,3.5)-- (3.5,4.);
\draw [line width=0.5pt, dash pattern=on 2pt off 2pt] (2.5,4.)-- (2.5,4.5);
\begin{scriptsize}
\draw [fill=black] (1.7,2.5) circle (1pt);
\draw [fill=black] (2.,2.) circle (1pt);
\draw [fill=black] (2.3,2.5) circle (1pt);
\draw [fill=black] (1.7,3.) circle (1pt);
\draw [fill=black] (1.5,3.5) circle (1pt);
\draw [fill=black] (2.,3.5) circle (1pt);
\draw [fill=black] (2.5,3.5) circle (1pt);
\draw [fill=black] (3.,3.5) circle (1pt);
\draw [fill=black] (2.,4.) circle (1pt);
\draw [fill=black] (2.5,4.) circle (1pt);
\draw [fill=black] (3.,4.) circle (1pt);
\draw [fill=black] (3.5,4.) circle (1pt);
\draw [fill=black] (2.5,4.5) circle (1pt);
\end{scriptsize}
\draw(2.4,2.1) node {\tiny $d$};
\draw(2.3,2.8) node {\tiny $c$};
\draw(1.5,2.6) node {\tiny $b$};
\draw(1.5,3.1) node {\tiny $a$};
\draw(1.3,3.6) node {\tiny $e$};
\draw(1.8,3.6) node {\tiny $f$};
\draw(2.3,3.6) node {\tiny $g$};
\draw(3.2,3.4) node {\tiny $h$};
\draw(1.8,4.1) node {\tiny $i$};
\draw(2.3,4.1) node {\tiny $j$};
\draw(2.8,4.1) node {\tiny $k$};
\draw(3.7,4.1) node {\tiny $l$};
\draw(2.3,4.6) node {\tiny $m$};
\end{tikzpicture}\end{gather*}
with $a,\dots,m\in \D$, then
\begin{gather*}
T_0^{(a)}=\tdtroisun{$d$}{$c$}{$b$}{10}{2},\qquad\qquad
\big\{T_1^{(a)},\dots,T_n^{(a)}\big\}=\begin{Bmatrix}
\tdun{$e$},\tddeux{$f$}{$i$}{10},
\tdtroisdeux{$g$}{$j$}{$m$}{2}{2},\tdtroisun{$h$}{$l$}{$k$}{10}{2}
\end{Bmatrix}\!.
\end{gather*}
\end{Example}

\begin{Proposition}
For any nonempty finite set $A$, we denote by $\P_\T(A)$ the vector space gene\-ra\-ted by $\bfT_\T(A)$.
We define a composition $\circ$ on $\P_\T$ in the following way: for any $T\in \bfT_\T(A)$,
$T'\in \bfT_\T(B)$ and $a\in A$,
\begin{gather*}
T\circ_a T'=\sum_{v_1,\dots,v_{n_a}\in V(T')}
\big(\cdots\big (\big(T_0^{(a)}\bullet_{a'}^{(t_0)} T'\big)\bullet_{v_1}^{(t_1)} T_1^{(a)}\big)\cdots\big)
\bullet_{v_{n_a}}^{(t_{n_a})} T_{n_a}^{(a)},\end{gather*}
where $a'$ is the direct ascendant of $a$ in $T$ and $t_0$ is the type of the edge between $a'$ and $a$.
If $a$ is the root of $T$, by convention $T_0^{(a)}\bullet_{a'}^{(t_0)} T'=T'$.
With this composition, $\P_\T$ is an operad in the category of species.
\end{Proposition}

\begin{proof}
Note that the tree
$\big(\cdots \big(\big(T_0^{(a)}\bullet_\lambda^{(t_0)} T'\big)\bullet_{v_1}^{(t_1)} T_1^{(a)}\big)\cdots\big)
\bullet_{v_{n_a}}^{(t_{n_a}^{(a)})} T_{n_a}$,
which is shortly denoted by $T\bullet_\lambda^{(v)} T'$, is obtained in the following process:
\begin{enumerate}\itemsep=0pt
\item Delete the branches $T_1^{(a)},\dots,T_{n_a}^{(a)}$ coming from $a$ in $T$.
One obtains a tree $T''$, and $a$ is a~leaf of $T''$.
\item Identify $a\in V(T'')$ with the root of $T'$. One obtains a tree $T'''$.
\item Graft $T_1^{(a)}$ on the vertex $v_1$ of $T'''$v with an edge of type $t_1$, $\dots$,
graft $T_{n_a}^{(a)}$ on the vertex~$v_{n_a}$ of $T'''$ with an edge of type $t_n$.
The obtained tree is $T\bullet_\lambda^{(v)} T'$.
\end{enumerate}
This obviously does not depend on the choice of the indexation of $T_1^{(a)},\dots,T_{n_a}^{(a)}$.

Let $T\in \bfT_\T(A)$, $T'\in \bfT_\T(B)$, $T''\in \bfT_\T(C)$.
\begin{itemize}\itemsep=0pt
\item If $a',a''\in A$, with $a'\neq a''$, then
\begin{align*}
(T\circ_{a'} T')\circ_{a''}T''&=\sum_{v'\in V(T')^{n_{a'}},v''\in V(T'')^{n_{a''}}}
\big(T\bullet_{a'}^{(v')}T'\big)\bullet_{a''}^{(v'')} T''
\\
&=\sum_{v'\in V(T')^{n_{a'}},v''\in V(T'')^{n_{a''}}}\big(T\bullet_{a''}^{(v'')}T''\big)\bullet_{a'}^{(v')} T'
\\
&=(T\circ_{a''} T'')\circ_{a'}T'.
\end{align*}
\item If $a'\in A$ and $b''\in B$, then
\begin{align*}
(T\circ_{a'}T')\circ_{b''}T''&=\sum_{v'\in V(T')^{n_{a'}},v''\in V(T'')^{n_{b''}}}
\big(T\bullet_{a'}^{(v')}T'\big)\bullet_{b''}^{(v'')} T''
\\
&=\sum_{v'\in V(T')^{n_{a'}},v''\in V(T'')^{n_{b''}}}T\bullet_{a'}^{(v')}\big(T'\bullet_{b''}^{(v'')} T''\big)
\\
&=T\circ_{a'}(T'\circ_{b''}T'').
\end{align*}
\end{itemize}
Moreover, $\tdun{$a$}\bullet_\lambda T=T$ for any tree $T$, and if $a\in V(T)$, $T\bullet_\lambda \tdun{$a$}T$.
So $\P_\bfT$ is indeed an operad in the category of species. \end{proof}

Consequently, the family $(\P_\T(n))_{n\geq 0}$ is an operad in the category of vector spaces, which we denote by $\P_\T$.

\begin{Example}
\begin{gather*}
\tddeux{$1$}{$2$}{10}\circ_1\tddeux{$1$}{$2$}{2}=\tdtroisun{$1$}{$3$}{$2$}{2}{10}
+\tdtroisdeux{$1$}{$2$}{$3$}{2}{10},\qquad\qquad
\tddeux{$1$}{$2$}{2}\circ_2\tddeux{$1$}{$2$}{10}=\tdtroisdeux{$1$}{$2$}{$3$}{2}{10}.
\end{gather*}
\end{Example}

\begin{Remark}
Another operad on typed trees is introduced in \cite{Dotsenko2}.
It is a typed version of the operad of nonassociative, permutative operad of \cite{Livernet},
and is different from ours.
\end{Remark}

In the non-typed case, this theorem is proved in \cite{ChapotonLivernet}:

\begin{Theorem} \label{theo11}
The operad of $\T$-multiple pre-Lie algebras is isomorphic to $\P_\T$,
via the isomorphism $\Phi$ sending, for any $t\in \T$, $\bullet_t$ to the tree
$\tddeux{$1$}{$2$}{10}$, where the edge is of type $t$.
\end{Theorem}

\begin{proof} The operad of $\T$-multiple pre-Lie algebras is generated by the binary elements $\bullet_t$,
$t\in \T$, with the relations
\begin{gather*}
\forall t,t'\in \T,\qquad
\bullet_{t'}\circ_2 \bullet_t-\bullet_t\circ_1\bullet_{t'}
=(\bullet_t\circ_2\bullet_{t'}-\bullet_{t'}\circ_1 \bullet_t)^{(23)},
\end{gather*}
where in the right side we used the action of the symmetric group $\mathfrak{S}_3$ on $\T(3)$,
and more specifically the action of the transposition $(23)$.
Firstly, if $t$ and $t'$ are elements of $\T$, symbolized by$\typ{10}$ and $\typ{2}$\!, by the preceding example:
\begin{gather*}
\tddeux{$1$}{$2$}{10}\circ_1\tddeux{$1$}{$2$}{2}-\tddeux{$1$}{$2$}{2}\circ_2\tddeux{$1$}{$2$}{10}
=\tdtroisun{$1$}{$3$}{$2$}{2}{10}
=\Big(\tdtroisun{$1$}{$3$}{$2$}{10}{2}\Big)^{(23)}
=\Big(\tddeux{$1$}{$2$}{2}\circ_1\tddeux{$1$}{$2$}{10}-\tddeux{$1$}{$2$}{10}\circ_2\tddeux{$1$}{$2$}{2}
\Big)^{(23)}.
\end{gather*}
So the morphism $\phi$ exists. Let us prove that it is surjective:
let $T\in \bfT_\T(n)$, we show that it belongs to $\im(\Phi)$ by induction on $n$. It is obvious if $n=1$ or $n=2$.
Let us assume the result at all ranks $<n$. Up to a reindexation we assume that
\begin{gather*}
T=B_1(T_1\partial_{t_1}\cdots T_k\partial_{t_k}),
\end{gather*}
 where for any $1\leq i<j\leq k$, if $x\in V(T_i)$ and $y\in V(T_j)$, then $x<y$.
We denote by $T'_i$ the standardization of $T_i$. By the induction hypothesis on $n$, $T'_i\in \im(\Phi)$ for all $i$.
We proceed by induction on $k$. The type $t_k$ will be represented in red. If $k=1$, then
\begin{gather*}
T=\tddeux{$1$}{$2$}{10}\circ_2 T_1\in \im(\Phi).
\end{gather*}
Let us assume the result at rank $k-1$. We put $T'=B_1(T_1\partial_{t_1}\cdots T_{k-1}\partial_{t_{k-1}})$.
By the induction hypothesis on $n$, $T'\in \im(\Phi)$. Then
\begin{gather*}
\tddeux{$1$}{$2$}{10}\circ_1 T'=T+x,
\end{gather*}
where $x$ is a sum of trees with $n$ vertices, such that the fertility of the root is $k-1$.
Hence, $x\in \im(\Phi)$, so $T\in \im(\Phi)$.

Let $\D$ be a set. The morphism $\phi$ implies that the free $\P_\bfT$-algebra generated by $\D$,
that is to say $\g_{\D,\T}$, inherits
a $\T$-multiple pre-Lie algebra structure defined by
\begin{gather*}
\forall x,y\in \g_{\D,\T},\quad \forall \typ{2}\in \T,\qquad
x\circ_{\typ{2}} y=\tddeux{$1$}{$2$}{2}\cdot(x\otimes y),
\end{gather*}
where $\cdot$ is the $\P_\T$-algebra structure of $\g_{\D,\T}$.
For any trees $T$, $T'$ in $\bfT_{\D,\T}$, by definition of the operadic composition of $\P_\T$,
\begin{gather*}
T\circ_t T'=\sum_{v\in V(T)} T\bullet_t^{(v)} T',
\end{gather*}
so $\circ_t=\bullet_t$ for any $t$. As $(\g_{\D,\T},(\bullet_t)_{t\in \T})$
is the free $\T$-multiple pre-Lie algebra generated by~$\D$,~$\Phi$~is an isomorphism.
\end{proof}
